\documentclass[12pt]{article}

\usepackage{amssymb,amsmath,amsthm}

% gaya teorema (semuanya memakai pencacah yang sama)
\theoremstyle{plain}
\newtheorem{theorem}{Teorema}[section]
\newtheorem{lemma}[theorem]{Lema}
\newtheorem{proposition}[theorem]{Proposisi}
\newtheorem{corollary}[theorem]{Korolari}
\newtheorem{conjecture}[theorem]{Konjektur}
\newtheorem{exercise}[theorem]{Latihan}
\theoremstyle{definition}
\newtheorem{definition}[theorem]{Definisi}
\newtheorem{example}[theorem]{Contoh}
\theoremstyle{remark}
\newtheorem{remark}[theorem]{Catatan}

% beberapa pengaturan penomoran
\numberwithin{equation}{section}

\usepackage{fancyhdr}
\usepackage{lastpage}
\setlength{\headheight}{15.2pt}
\renewcommand{\footrulewidth}{0.4pt}% nilai bakunya 0pt
\setlength{\footskip}{30pt}
\pagestyle{fancy}

\makeatletter
\let\ps@plain\ps@fancy 
\makeatother

\lhead{MATH 742}
\chead{Modul dan Teori Wedderburn}
\rhead{Musim Semi 2023}
\lfoot{Terakhir Direvisi: \today}
\cfoot{}
\rfoot{\thepage\ dari \pageref{LastPage} }

\begin{document}

%%%%%%%%%%%%%%%%%%
%%%%%%%%%%%%%%%%%%
\title{Modul dan Teori Wedderburn}
\author{Alexander Duncan}

\maketitle

Teori representasi modern biasanya dirumuskan dalam bahasa
gelanggang/aljabar nonkomutatif dan modul-modulnya.
Sampai sekarang saya menghindari bahasa ini agar dapat berfokus pada dasar-dasarnya.
Namun, sudut pandang teori modul jauh lebih serbaguna --- bahkan
untuk grup hingga atas $\mathbb{C}$.
Sudut pandang ini jauh lebih mendesak diperlukan ketika membahas teori
representasi atas medan-medan lain. 

Objek utamanya adalah \emph{gelanggang grup} $\mathbb{Z}G$
atau \emph{aljabar grup} $kG$. Kita akan melihat bahwa sebagian besar bahasa
teori representasi grup hingga dapat dinyatakan secara efisien
sebagai kasus khusus teori modul atas $kG$.
Kemudian kita akan memakai bahasa tersebut untuk melangkah lebih jauh.

Karena teori modul merupakan bagian dari rangkaian ujian kualifikasi aljabar baku,
saya akan sering menghilangkan bukti, bahkan banyak definisi, dengan
asumsi bahwa pembaca sudah pernah menjumpainya.
Namun, mata kuliah persiapan ujian kualifikasi biasanya menekankan gelanggang
\emph{komutatif}, sehingga beberapa pengingat tetap diperlukan.
Saya menyarankan pembaca mencari buku teks aljabar tingkat pascasarjana yang
membahas modul dengan baik, seperti \cite[\S{10}]{DF} atau \cite[\S{III}]{Lang},
untuk dijadikan rujukan jika ada hal yang belum dikenal.

“Dasar-dasar” teori representasi berbasis modul terdapat dalam
hampir setiap teks yang telah saya rujuk, sering kali pada tahap yang sangat awal; 
lihat \cite{DF}, \cite{Etingof}, \cite{FultonHarris}, \cite{Lang}, dan sebagainya.
Buku Serre \cite{Serre} mula-mula tidak memakai sudut pandang teori
modul, tetapi kemudian mendadak mengubah haluan pada bab 6 dan
mengasumsikan pembaca sudah mengenal teori Wedderburn (yang dikembangkan di bawah).
Saya juga akan banyak mengambil materi dari \cite[\S{12,13}]{AlperinBell},
\cite{AlperinLRT}, dan \cite{CurtisReiner}.

\subsection*{Konvensi}

Sepanjang pembahasan, kita mengasumsikan semua gelanggang dan aljabar mempunyai satuan
(memiliki unsur identitas), tetapi \emph{tidak harus} komutatif.
Kita mencadangkan $k$ untuk gelanggang komutatif, yang biasanya berupa
medan atau $\mathbb{Z}$.
Ketika gelanggang $R$ komutatif, kita biasanya mengabaikan
pembedaan antara modul kiri dan kanan atas $R$, lalu
menyebut keduanya modul atas $R$.

%%%%%%%%%%%%%%%%%%
%%%%%%%%%%%%%%%%%%

\section{Pemanasan: Aljabar Matriks atas Medan}

Dalam bagian ini, kita meninjau kasus khusus aljabar nonkomutatif
atas $k$, yaitu $\operatorname{M}_n(k)$, dengan $k$ suatu medan.
Banyak kekhasan gelanggang nonkomutatif sudah tampak bahkan dalam
konteks ini.

\begin{example}
Ruang $k^n$ dari vektor-vektor kolom merupakan modul kiri atas $\operatorname{M}_n(k)$.
Ruang $(k^n)^T$ dari vektor-vektor baris merupakan modul kanan atas $\operatorname{M}_n(k)$.
Secara lebih umum, untuk bilangan bulat $m$, ruang matriks berukuran $n \times m$
$\operatorname{M}_{n \times m}(k)$ merupakan modul kiri atas $\operatorname{M}_n(k)$,
dan ruang matriks berukuran $m \times n$
$\operatorname{M}_{m \times n}(k)$ merupakan modul kanan atas $\operatorname{M}_n(k)$.
\end{example}

Sebentar lagi kita akan melihat bahwa contoh di atas telah
menampilkan semua kelas isomorfisme modul yang terbangkitkan hingga,
baik kiri maupun kanan, atas $\operatorname{M}_{m \times n}(k)$.

\begin{definition}
Misalkan $R$ suatu gelanggang dan $M$ suatu modul (kiri atau kanan) atas $R$.
Kita menyebut $M$ \emph{sederhana} jika $M\ne 0$ tidak mempunyai submodul selain
$0$ dan $M$ sendiri.
Kita menyebut $M$ \emph{tak terdekomposisi} jika $M$ tidak dapat ditulis sebagai jumlah langsung
submodul-submodul tak nol.
Kita menyebut $M$ \emph{semisederhana} jika $M$ merupakan jumlah langsung
submodul-submodul sederhana.
\end{definition}

Dari contoh di atas, kita melihat bahwa $k^n$ bersifat sederhana sekaligus
tak terdekomposisi sebagai modul kiri atas $\operatorname{M}_n(k)$.
Kita mempunyai dekomposisi
\[
\operatorname{M}_{n \times m}(k) \cong (k^n)^{\oplus m}
\]
sebagai modul kiri atas $\operatorname{M}_n(k)$,
sehingga $\operatorname{M}_{n \times m}(k)$ tidak sederhana maupun
tak terdekomposisi; namun, modul itu semisederhana.

\subsection{Ideal}

Diberikan subruang $V \subseteq k^n$, misalkan
$I$ adalah subhimpunan $\operatorname{M}_n(k)$
yang setiap barisnya diambil dari $V$,
dan misalkan $J$ adalah subhimpunan $\operatorname{M}_n(k)$
yang setiap kolomnya diambil dari $V$.
Perhatikan bahwa perkalian kiri pada $I$ memetakan setiap baris ke suatu
kombinasi linear unsur-unsur $V$, sehingga kita menyimpulkan bahwa $I$ merupakan
ideal kiri. Serupa dengan itu, $J$ merupakan ideal kanan.

Prosedur ini dapat dibalik.
Diberikan ideal kiri $I$ dari $\operatorname{M}_n(k)$,
kita dapat merekonstruksi subruang $V$
sebagai himpunan baris yang muncul di antara unsur-unsur $I$.
Demikian pula untuk ideal kanan. Kita menyimpulkan:

\begin{proposition}
Ideal kiri (berturut-turut, ideal kanan) dari $\operatorname{M}_n(k)$
berkorespondensi bijektif secara kanonik dengan subruang-subruang vektor dari $k^n$.
\end{proposition}

\begin{example}
Menurut saya, korespondensi-korespondensi tersebut entah bagaimana lebih mudah dipahami
melalui contoh konkret. Misalkan $V$ suatu subruang,
$I$ ideal kiri yang bersesuaian, dan $J$ ideal kanan yang bersesuaian.
Nilai $a,b,c,d,e,f$ menyatakan unsur-unsur sebarang dari $k$.
Dalam $\operatorname{M}_2(k)$, kita mempunyai contoh:
\[
V = \operatorname{span}_k \left\{
\begin{pmatrix} 1 \\ 0 \end{pmatrix} \right\}
\quad
I = \left\{ \begin{pmatrix} a & 0 \\ b & 0 \end{pmatrix} \right\} 
\quad
J = \left\{ \begin{pmatrix} a & b \\ 0 & 0 \end{pmatrix} \right\}
\]
Dalam $\operatorname{M}_3(k)$, kita mempunyai
\[
V = \operatorname{span}_k \left\{
\begin{pmatrix} 1 \\ 1 \\ 0 \end{pmatrix},
\begin{pmatrix} 0 \\ 1 \\ 1 \end{pmatrix} \right\}
\quad
J = \left\{ \begin{pmatrix}
a & b & c \\
a+d & b+e & c+f \\
d & e & f
\end{pmatrix} \right\}
\]
dan $I$ hanyalah transposnya.
\end{example}

\begin{remark}
Bagi pembaca yang mengenal sedikit geometri aljabar,
himpunan subruang berdimensi tetap $r$ dalam $k^n$ adalah
\emph{Grassmannian} $\operatorname{Gr}(n,r)$.
Khususnya, himpunan subruang berdimensi $1$ adalah ruang projektif
$\mathbb{P}^{n-1}$.
Korespondensi di atas memungkinkan kita mendefinisikan Grassmannian dan
ruang projektif sebagai himpunan ideal kiri dari $\operatorname{M}_n(k)$.
Pergeseran sudut pandang ini berguna ketika kita mendefinisikan objek yang
lebih umum, seperti varietas Severi--Brauer, ketika aljabarnya bukan aljabar
matriks.
\end{remark}

Dengan memilih basis $v_1,\ldots v_m$ bagi subruang $V$, perhatikan matriks
\[
M_V := \begin{pmatrix}
\vdots & \cdots & \vdots & \vdots & \cdots & \vdots\\
v_1 & \cdots & v_n & 0 & \cdots & 0\\
\vdots & \cdots & \vdots & \vdots & \cdots & \vdots\\
\end{pmatrix} 
\]
yang $m$ kolom pertamanya merupakan vektor-vektor basis $v_1,\ldots,v_m$ dan
kolom sisanya bernilai $0$.
Kita mengamati bahwa ideal kanan $J$ yang bersesuaian dibangkitkan oleh $M_V$.
Serupa dengan itu, $I$ dibangkitkan oleh $M_V^T$.

Ingat bahwa suatu ideal disebut \emph{utama}
jika dibangkitkan oleh satu unsur.
Dengan demikian, kita memperoleh:

\begin{corollary}
Setiap ideal kiri (berturut-turut, ideal kanan) dari $\operatorname{M}_n(k)$ bersifat utama.
\end{corollary}

\begin{definition}
Kita menyebut gelanggang $R$ \emph{sederhana} jika dua-sisi idealnya hanya $0$ dan
$R$ sendiri.
Kita menyebut gelanggang $R$ \emph{semisederhana} jika merupakan hasil kali
gelanggang-gelanggang sederhana.
\end{definition}

Dari deskripsi eksplisit ideal kiri dan kanan
dari $\operatorname{M}_n(k)$, kita melihat bahwa:

\begin{proposition}
$\operatorname{M}_n(k)$ merupakan gelanggang sederhana.
\end{proposition}

Berhati-hatilah terhadap kemungkinan ketaksaan pada kata
“sederhana”.
Ingat bahwa $\operatorname{M}_n(k)$ juga dapat dipandang sebagai modul kiri
(atau kanan) atas dirinya sendiri.
Sebagai modul kiri,
$\operatorname{M}_n(k)$ biasanya \emph{tidak} sederhana karena mempunyai banyak
submodul (ideal kiri).

Kelak kita akan melihat bahwa aljabar berdimensi hingga semisederhana sebagai
gelanggang jika dan hanya jika semisederhana sebagai modul (kiri atau kanan),
tetapi hal ini tidak jelas: suku-suku sederhana dalam kedua
konteks tersebut tidak sama!

\subsection{Idempoten}

\begin{definition}
Misalkan $R$ suatu gelanggang. Unsur $e$ disebut \emph{idempoten} jika $e^2=e$.
Dua idempoten $e_1,e_2$ disebut \emph{ortogonal} jika $e_1e_2=e_2e_1=0$.
Sebuah \emph{idempoten primitif} adalah idempoten tak nol
yang bukan jumlah dua idempoten ortogonal tak nol.
\end{definition}

Idempoten-idempoten $\operatorname{M}_n(k)$ tepat merupakan (matriks yang
merepresentasikan) proyeksi $\pi : k^n \to k^n$.
Diberikan idempoten $e \in \operatorname{M}_n(k)$ dengan proyeksi
yang bersesuaian $\pi_e$, kita memperoleh dekomposisi jumlah langsung
\[
k^n = \ker(\pi_e) \oplus \operatorname{im}(\pi_e)
\ .
\]
Perhatikan bahwa idempoten $1-e$ ortogonal terhadap $e$ dan bahwa
\[
\ker(\pi_{1-e}) = \operatorname{im}(\pi_{e})
\textrm{ dan }
\operatorname{im}(\pi_{1-e}) = \ker(\pi_e).
\]
Jadi, kita mempunyai dekomposisi ruang vektor
\[
k^n = e(k^n) + (1-e)(k^n)
\]
dan kita ingatkan bahwa $e(k^n)$ serta $(1-e)(k^n)$ \emph{bukan}
modul atas $\operatorname{M}_n(k)$.

Secara lebih umum, himpunan idempoten ortogonal $e_1,\ldots,e_r$
sedemikian sehingga $1=e_1+\ldots+e_r$ memberikan dekomposisi
\[
k^n = V_1 \oplus V_2 \oplus \cdots \oplus V_r
\]
dengan $V_i = \operatorname{im}(\pi_i) = e_i(k^n)$ untuk $1 \le i \le r$.

Kita melihat bahwa idempoten primitif dari $\operatorname{M}_n(k)$
bersesuaian dengan proyeksi berperingkat $1$.
Selain itu, suatu himpunan maksimal $e_1,\ldots,e_r$ dari idempoten
primitif ortogonal sama artinya dengan pilihan basis bagi $k^n$.

\begin{example}
Perhatikan aljabar matriks $\operatorname{M}_2(\mathbb{C})$
dan tetapkan
\[
a = \begin{pmatrix} 1 & 0 \\ 0 & 0 \end{pmatrix},
b = \begin{pmatrix} 0 & 0 \\ 0 & 1 \end{pmatrix},
c = \begin{pmatrix} \frac{1}{2} & \frac{1}{2} \\ \frac{1}{2} & \frac{1}{2} \end{pmatrix},
\textrm{ dan }
d = \begin{pmatrix} \frac{1}{2} & -\frac{1}{2} \\ -\frac{1}{2} & \frac{1}{2} \end{pmatrix}.
\]
Masing-masing $a,b,c,d$ merupakan idempoten primitif.
Kita melihat bahwa $a,b$ ortogonal dan $c,d$ ortogonal.
Namun, $a,c$ tidak ortogonal.
\end{example}

Idempoten merupakan bahasa yang praktis untuk memahami
dekomposisi ideal.
Jika $V$ adalah subruang yang bersesuaian dengan ideal kiri $I$,
maka $I = \operatorname{M}_n(k)e$, dengan $e$ idempoten yang
bersesuaian dengan proyeksi $k^n \to V$.

Secara lebih umum, jika $e_1,\ldots,e_r$ adalah himpunan idempoten
ortogonal yang memenuhi $e_1+\cdots+e_r=1$, maka
\[
\operatorname{M}_n(k) = \operatorname{M}_n(k)e_1
\oplus \cdots \oplus \operatorname{M}_n(k)e_r
\]
sebagai modul kiri atas $\operatorname{M}_n(k)$, dan
\[
\operatorname{M}_n(k) = e_1\operatorname{M}_n(k)
\oplus \cdots \oplus e_r\operatorname{M}_n(k)
\]
sebagai modul kanan atas $\operatorname{M}_n(k)$.
Selain itu, kita mempunyai dekomposisi
\[
\operatorname{M}_n(k) = \bigoplus_{1\le i,j \le r}
e_i\operatorname{M}_n(k)e_j
\]
sebagai ruang vektor atas $k$, yang dapat dipandang sebagai sejenis dekomposisi
matriks blok.

\begin{definition}
Misalkan $R$ suatu gelanggang.
\emph{Pusat} dari $R$ adalah subhimpunan
\[
Z(R) := \{ z \in R \mid rz=zr \textrm{ untuk semua } r \in R \},
\]
yang merupakan subgelanggang komutatif dari $R$.
\end{definition}

Perhatikan bahwa $Z(\operatorname{M}_n(k)) \cong k$ hanyalah himpunan matriks
skalar.

\begin{definition}
Idempoten $e$ disebut \emph{sentral} jika $e$ termuat dalam pusat
$Z(R)$ dari $R$.
Sebuah \emph{idempoten sentral primitif} adalah idempoten sentral tak nol
yang bukan jumlah dua idempoten sentral ortogonal tak nol.
\end{definition}

Peringatan: idempoten sentral primitif selalu sentral,
tetapi belum tentu primitif!
Kadang-kadang istilah “primitif secara sentral” dipakai
sebagai ganti “sentral primitif”, yang terasa lebih atau kurang membingungkan,
bergantung pada bagaimana suasana hati Anda pagi itu.

Aljabar matriks $\operatorname{M}_n(k)$ mempunyai tepat dua idempoten
sentral: matriks nol dan matriks identitas.
Hanya matriks identitas yang merupakan idempoten sentral primitif.
(Identitas \emph{bukan} idempoten primitif kecuali $n=1$.)

Idempoten sentral tidak begitu menarik bagi gelanggang matriks, tetapi
menarik bagi hasil kali gelanggang.

\begin{exercise}
Misalkan $R$ suatu gelanggang dan $e$ suatu idempoten sentral.
Maka $eRe$ merupakan gelanggang terhadap penjumlahan dan perkalian yang sama,
dengan $e$ sebagai identitas.
(Perhatikan bahwa $eRe$ menjadi subgelanggang hanya ketika $eRe$ memuat
identitas, yang hanya terjadi jika $e=1$.)
\end{exercise}

\begin{exercise}
Misalkan $R$ suatu gelanggang. Sebuah dekomposisi
\[
R= R_1 \times \cdots \times R_n
\]
dengan $R_1,\ldots,R_n$ gelanggang-gelanggang berkorespondensi dengan himpunan
idempoten sentral ortogonal $e_1,\ldots,e_n$ sedemikian
sehingga $R_i \cong e_iRe_i$ untuk setiap $1 \le i \le n$.
\end{exercise}

\begin{exercise}
Misalkan $R$ suatu gelanggang dan $e$ suatu idempoten sentral primitif.
Maka $eRe$ tidak dapat ditulis sebagai hasil kali dua gelanggang tak nol.
\end{exercise}

%%%%%%%%%%%%%%%%%%
%%%%%%%%%%%%%%%%%%


\section{Gelanggang Grup}

\begin{definition}
Misalkan $G$ suatu grup hingga.
\emph{Gelanggang grup} $\mathbb{Z}G$ adalah grup abelian bebas
pada $G$ dengan perkalian yang diberikan oleh
\[
\left( \sum_{g \in G} a_g g \right) \cdot
\left( \sum_{g \in G} b_g g \right)
:= \sum_{g \in G} \sum_{h \in G} a_gb_h gh
\]
untuk bilangan bulat $\{a_g\}_{g\in G}$ dan $\{b_g\}_{g\in G}$.
Secara lebih umum, diberikan gelanggang komutatif $k$,
\emph{aljabar grup} $kG$ adalah modul bebas atas $k$ pada $G$
dengan perkalian seperti di atas.
\end{definition}

Langsung dapat diperiksa bahwa $\mathbb{Z}G$ merupakan gelanggang nonkomutatif dengan
identitas $1$ yang bersesuaian dengan unsur basis dari identitas dalam $G$.
Serupa dengan itu, $kG$ merupakan $k$-aljabar dengan identitas $1$.
Karena $1 \in G$ bersesuaian dengan $1 \in \mathbb{Z}G$,
mengidentifikasi basis $\mathbb{Z}G$ dengan unsur-unsur
$G$ pada umumnya tidak menimbulkan masalah jika $G$ ditulis secara multiplikatif.

\begin{example}
Misalkan $G=D_6=\langle s,r \mid s^2, r^3, (sr)^2 \rangle$
adalah grup dihedral berorde $6$.
Kita mempunyai perhitungan berikut dalam $\mathbb{Z}G$:
\begin{align*}
&(3 + 2s + 5r)(r-sr + 5r^2)\\
=& 3r -3sr +15r^2 + 2sr -2r+10sr^2 +5r^2 -5s+25\\
=& 25 -5s + r + 20r^2 -sr +10sr^2
\end{align*}
\end{example}

\begin{example}
Misalkan $G = \langle r \mid r^n \rangle$ adalah grup siklik berorde $n$.
Jadi, $G$ mempunyai basis $\{1,r,r^2,\ldots, r^{n-1} \}$
dengan $r^i\cdot r^j=r^{i+j}$ tunduk pada relasi $r^n=1$.
Dengan kata lain, $\mathbb{Z}G$ isomorfik dengan gelanggang
$\mathbb{Z}[x]/(x^n-1)$.
\end{example}

Perhatikan bahwa memakai notasi aditif $\mathbb{Z}/n\mathbb{Z}$
untuk grup siklik berorde $n$ dapat menimbulkan ketaksaan
yang berbahaya dalam contoh di atas.
Karena alasan ini, kita biasanya memakai “notasi eksponensial”
untuk gelanggang grup dari grup aditif.
Jika $a+b=c$ dalam suatu grup aditif, maka
kita tulis $x^ax^b=x^{a+b}=x^c$ dalam gelanggang grupnya.

\begin{exercise}
Buktikan bahwa $\mathbb{Z}G$ komutatif jika dan hanya jika $G$ abelian.
\end{exercise}

Berguna untuk menuliskan koefisien hasil kali sebagai rumus dalam
faktor-faktornya. Secara khusus, jika
\[
\left( \sum_{g \in G} a_g g \right) \cdot
\left( \sum_{g \in G} b_g g \right)
= \sum_{g \in G} c_g g
\]
maka
\[
c_g = \sum_{h \in G} a_hb_{h^{-1}g}
\]
untuk setiap $g \in G$.

Karena $G$ hingga, deskripsi ekuivalen bagi $\mathbb{Z}G$ adalah himpunan
$\operatorname{Hom}_{\mathrm{Set}}(G,\mathbb{Z})$
fungsi-fungsi dari $G$ ke $\mathbb{Z}$.
Namun, struktur gelanggangnya \emph{bukan} struktur yang “sudah jelas”.
Struktur aditifnya memang berupa penjumlahan titik demi titik,
tetapi kita memakai \emph{hasil kali konvolusi}
dari dua fungsi $f_1,f_2$ yang didefinisikan melalui
\[
(f_1 \ast f_2)(g) = \sum_{ij=g} f_1(i)f_2(j)
\]
untuk $g \in G$, dengan indeks penjumlahan $i,j$
bervariasi atas semua pasangan unsur $G$ sedemikian sehingga $ij=g$.
Isomorfisme dengan $\mathbb{Z}G$ diperoleh dengan mengidentifikasi setiap unsur basis
$g \in \mathbb{Z}G$ dengan fungsi karakteristik $\chi_g$
dalam $\operatorname{Hom}_{\mathrm{Set}}(G,\mathbb{Z})$
yang memenuhi $\chi_g(h) = \delta_{gh}$ untuk semua $h \in G$.

\subsection{Modul dari Gelanggang Grup}

Kaitan gelanggang grup dengan teori grup tampak dari
proposisi berikut:

\begin{proposition} \label{prop:equivalence}
Misalkan $G$ grup hingga dan $k$ medan.
Struktur representasi $G$ atas medan $k$
ekuivalen dengan struktur modul kiri atas $kG$.
Representasi berdimensi hingga bersesuaian dengan
modul yang terbangkitkan hingga.
Selain itu, transformasi linear yang $G$-ekuivarian tepat merupakan
homomorfisme modul atas $kG$.
\end{proposition}

\begin{proof}
Misalkan $M$ suatu modul kiri atas $kG$.
Khususnya, $M$ merupakan ruang vektor atas $k$.
Kita mengonstruksi representasi linear $\rho : G \to \operatorname{GL}(M)$
melalui $\rho_g(m)=gm$ untuk $g \in G$ dan $m\in M$.
Sebaliknya, diberikan representasi linear $\sigma : G \to
\operatorname{GL}(V)$, kita lengkapi $V$ dengan struktur modul kiri atas $kG$ melalui
\[
\left( \sum_{g \in G} a_g g \right) v := \sum_{g \in G} a_g \sigma_g(v)
\]
untuk $\{a_g\}$ dari $k$ dan $v \in V$.
Kedua operasi ini saling invers.

Jika $M$ merupakan modul kiri atas $kG$ yang terbangkitkan hingga,
misalnya oleh $\{s_1,\ldots,s_n\}$, maka $M$ sebagai ruang vektor atas $k$
direntang oleh $\{ gs_i \mid g \in G, 1 \le i \le n \}$.
Jadi, $\rho$ merupakan representasi berdimensi hingga.
Sebaliknya, jika $V$ berdimensi hingga, maka basis bagi $V$
sebagai ruang vektor atas $k$ terlebih lagi merupakan himpunan pembangkit bagi $V$ sebagai
modul atas $kG$.

Terakhir, kita mengamati bahwa transformasi linear yang $G$-ekuivarian
$f : (V,\rho) \to (W,\sigma)$ antara representasi tepat merupakan homomorfisme modul atas $kG$
menurut korespondensi di atas.
Memang,
\begin{align*}
&f\left( \left( \sum_{g\in G} c_g g \right) v \right)\\
=& f\left( \sum_{g\in G} c_g \rho_g(v) \right)\\
=& \sum_{g\in G} c_g f\left( \rho_g(v) \right)\\
=& \sum_{g\in G} c_g \sigma_g\left(f(v)\right)\\
=& \left(\sum_{g\in G} c_g g\right) f(v)
\end{align*}
dengan menggunakan fakta bahwa $f$ linear dan $G$-ekuivarian.
\end{proof}

Kita langsung melihat bahwa subrepresentasi bersesuaian dengan submodul,
representasi hasil bagi bersesuaian dengan modul hasil bagi, dan jumlah langsung
bersesuaian dengan jumlah langsung.
Namun, meskipun $V \otimes_k W$ dan $\operatorname{Hom}_k(V,W)$ mempunyai
struktur modul atas $kG$ yang terinduksi, keduanya \emph{tidak} sama dengan
konstruksi teori modul $V \otimes_{kG} W$
dan $\operatorname{Hom}_{kG}(V,W)$ yang dibahas di bawah.
Selain itu, “modul trivial” adalah representasi nol,
\emph{bukan} “representasi trivial”.

\begin{example}
Setiap gelanggang sekaligus merupakan modul kiri dan kanan atas dirinya sendiri.
Aljabar grup $kG$ yang dipandang sebagai modul kiri atas $kG$ bersesuaian dengan
representasi reguler $G$.
\end{example}

Kehalusan penting lainnya di atas ialah bahwa dimensi suatu
representasi secara umum \emph{tidak} sama dengan banyak minimal
pembangkit modul atas $kG$ yang bersesuaian.

Modul kiri atas $kG$ tak terdekomposisi jika modul itu tak terdekomposisi sebagai
representasi.
Modul kiri atas $kG$ sederhana jika dan hanya jika tak tereduksi
sebagai representasi $G$.
Khususnya, representasi tak tereduksi dapat
dibangkitkan oleh satu unsur sebagai modul kiri atas $kG$ (meskipun kebalikannya mungkin
tidak benar).

\begin{theorem}[Teorema Maschke, dirumuskan kembali]
Jika $G$ hingga dengan orde yang relatif prima terhadap karakteristik medan $k$,
maka setiap modul kiri atas $kG$ yang terbangkitkan hingga bersifat semisederhana.
\end{theorem}

\subsection{Pusat Gelanggang Grup}

Pusat gelanggang grup sangat penting.
Perhatikan bahwa kita juga memiliki pengertian pusat $Z(G)$ dari grup $G$.
Memang, gelanggang grup dari pusat suatu grup termuat dalam
pusat gelanggang grup; dengan kata lain,
\[
kZ(G) \subseteq Z(kG).
\]
Namun, kesamaan hanya berlaku ketika $G=Z(G)$, sebagai akibat dari
hasil berikut.

\begin{proposition}
Pusat $Z(kG)$ dari aljabar grup $kG$ mempunyai basis
\[
\left\{
\sum_{g \in K} g\ \middle|\ K \in \mathcal{K}
\right\}
\]
dengan $\mathcal{K}$ himpunan kelas-kelas konjugasi dari $G$.
\end{proposition}

\begin{proof}
Misalkan $K$ suatu kelas konjugasi dalam $\mathcal{K}$.
Perhatikan bahwa $g \in K$ mengakibatkan
$hgh^{-1} \in K$.
Jadi, kita mengamati bahwa
\[
h\left(\sum_{g \in K} g\right)
=\left(\sum_{g \in K} g\right) h
\]
untuk semua $h \in H$. 
Sebaliknya, setiap unsur $x$ dari aljabar grup
yang memenuhi $hx=xh$ harus mempunyai koefisien yang sama
pada unsur-unsur basis dalam kelas konjugasi yang sama.
\end{proof}

\begin{example}
Misalkan $G=D_6=\langle s,r \mid s^2, r^3, (sr)^2 \rangle$
adalah grup dihedral berorde $6$.
Kita memperoleh bahwa
$Z(\mathbb{Z}G)$ mempunyai basis
\[
\{ 1, a = s+sr+sr^2, b = r+r^2 \}.
\]
Tabel perkaliannya ditentukan oleh beberapa perhitungan:
\begin{align*}
a^2 &= 3+3b\\
ab=ba &= 2a\\
b^2 &=2+b.
\end{align*}
Jadi, struktur pusatnya agak sukar dilihat dalam basis ini.
\end{example}

Ketika $k=\mathbb{C}$, pusatnya sangat mudah dideskripsikan.
Dalam Korolari~\ref{cor:CGcenter} di bawah, kita akan melihat bahwa
$Z(\mathbb{C}G) \cong \mathbb{C}^{|K|}$ sebagai $\mathbb{C}$-aljabar.
Namun, hal ini sama sekali tidak jelas dalam
basis yang diuraikan di atas!

\subsection{Dekomposisi Aljabar Grup Kompleks}

Misalkan $G$ suatu grup hingga.
Misalkan $W_1,\ldots,W_r$ adalah representasi
kompleks tak tereduksi yang berbeda dari $G$,
dan misalkan $n_1,\ldots, n_r$ dimensinya.

\begin{theorem}
Terdapat isomorfisme kanonik
\[
\mathbb{C}G \cong \bigoplus_{i=1}^r \operatorname{End}_{\mathbb{C}}(W_i)
\]
antara $\mathbb{C}$-aljabar.
\end{theorem}

\begin{proof}
Jika $V$ modul kiri atas $\mathbb{C}G$ dan $x$ unsur dari $kG$,
maka peta perkalian $v \mapsto xv$ merupakan endomorfisme
dari $V$.
Ini memberikan peta dari $\mathbb{C}G$ ke setiap
$\operatorname{End}_C(W_i)$ dan kita memperoleh
peta ke jumlah langsungnya.
Peta tersebut injektif karena aksi pada representasi reguler
setia.
Karena kedua aljabar berdimensi $n_1^2+\cdots+n_r^2$,
peta ini harus merupakan isomorfisme.
\end{proof}

Perhatikan bahwa untuk setiap ruang vektor kompleks $V$ berdimensi $n$, kita mempunyai
isomorfisme (nonkanonik)
$\operatorname{End}(V) \cong \operatorname{M}_n(\mathbb{C})$.
Jadi, teorema di atas dapat kita tulis kembali sebagai:  

\begin{corollary} \label{cor:fourier_as_matrices}
$\displaystyle
\mathbb{C}G \cong \bigoplus_{i=1}^r \operatorname{M}_{n_i}(\mathbb{C})$
\end{corollary}

Ingat bahwa pusat aljabar matriks
$\operatorname{M}_n(\mathbb{C})$ hanyalah subaljabar matriks
skalar, yang isomorfik dengan $\mathbb{C}$.
Berdasarkan Korolari~\ref{cor:fourier_as_matrices}, karena itu kita
memperoleh:

\begin{corollary} \label{cor:CGcenter}
$\displaystyle Z(\mathbb{C}G) \cong \mathbb{C}^r$
sebagai $\mathbb{C}$-aljabar.
\end{corollary}

Jadi, kita mempunyai dua basis bagi $Z(\mathbb{C}G)$: satu diindeks oleh kelas
konjugasi dan satu lagi diindeks oleh representasi tak tereduksi.
Tabel karakter $G$ tepat merupakan matriks perubahan basis antara
kedua basis ini.

\begin{example}
Misalkan $G = \mathbb{Z}/3\mathbb{Z}$ adalah grup siklik berorde $3$.
Perhatikan bahwa $kG \cong k[x]/(x^3-1)$ untuk setiap medan $k$.
Kita memperoleh $\mathbb{R}G \cong \mathbb{R} \oplus \mathbb{C}$,
sehingga korolari tersebut lebih halus atas medan yang tidak tertutup.
Selain itu, unsur $x-1$ dalam $\mathbb{F}_3G$ nilpoten,
sehingga keadaan berpotensi jauh lebih buruk ketika teorema Maschke
tidak berlaku.
\end{example}

\subsection{Idempoten Gelanggang Grup}

Sebelumnya dalam mata kuliah ini, kita telah melihat beberapa contoh proyeksi dalam
$\operatorname{End}_k^G(V)$ untuk berbagai representasi $(V,\rho)$.
Sebagai contoh, proyeksi $\pi : V \to V^G$ ke subruang
invarian mempunyai rumus
\[
\pi(v) = \frac{1}{|G|} \sum_{g \in G} \rho_g(v)
\]
untuk $v \in V$.
Contoh lain mencakup proyektor Young, penyimetris Young,
dan proyeksi ke komponen-komponen isotipik.
Semuanya lebih alami dipandang berasal dari unsur-unsur
aljabar grup.

Diberikan representasi $\rho : G \to \operatorname{GL}(V)$
dan unsur
\[
f = \sum_{g \in G} c_g g
\]
dalam aljabar grup $kG$,
definisikan $\widehat{f}(\rho) \in \operatorname{End}(V)$ melalui
\[
\widehat{f}(\rho)(v) = \sum_{g \in G} c_g \rho_g(v)
\]
untuk semua $v \in V$.
Jadi, proyeksi $\pi$ di atas hanyalah endomorfisme
$\widehat{f}(\rho)$,
dengan $f = |G|^{-1} \sum_{g \in G} g$.
Dalam praktiknya, kita cukup menulis $f$ sebagai ganti $\widehat{f}(\rho)$,
karena jarang ada risiko kebingungan.

Pertama, amati bahwa jika $p \in kG$ suatu idempoten,
maka $\widehat{p}(\rho)$ merupakan proyeksi untuk setiap representasi $\rho$.
(Kebalikannya tidak benar; perhatikan jumlah langsung dua salinan representasi
reguler dan proyeksi ke salah satu sukunya.)
Kedua, jika $p$ suatu idempoten primitif, maka proyeksi
yang bersesuaian $\pi$ pada representasi reguler bercitra suatu
subrepresentasi tak terdekomposisi.

Untuk gelanggang grup kompleks, kita bahkan dapat mengatakan lebih banyak.

\begin{proposition}
Misalkan $\chi_1,\ldots \chi_r$ adalah karakter kompleks tak tereduksi dari $G$.
Idempoten sentral primitif dari $\mathbb{C}G$ tepat merupakan
unsur-unsur
\[
e_i := \frac{\chi_i(1)}{|G|} \sum_{g\in G} \overline{\chi_i(g)} g
\]
untuk setiap $i=1,\ldots,r$.
Jika $(V,\rho)$ merupakan representasi $G$, maka
$\widehat{e_i}(\rho) : V \to V$ adalah proyeksi ke komponen
isotipik $V$ yang bersesuaian dengan $\chi_i$.
\end{proposition}

\begin{proof}
Dari isomorfisme
\[
\mathbb{C}G \cong \bigoplus_{i=1}^r \operatorname{End}(W_i)
\]
kita ingat bahwa pusatnya adalah $\mathbb{C}^r$,
yang bersesuaian dengan perkalian skalar dalam setiap
$\operatorname{End}(W_i)$.
Idempoten sentral dari $\mathbb{C}G$
bersesuaian dengan kasus ketika setiap perkalian skalar bernilai $1$ atau $0$.
Idempoten sentral primitif adalah idempoten yang menjadi identitas
dalam tepat satu $\operatorname{End}(W_i)$ dan trivial di tempat lain.
Rumus eksplisit untuk idempoten sentral primitif kini diperoleh
dari rumus eksplisit untuk proyeksi ke komponen-komponen
isotipik suatu representasi. 
\end{proof}

\section{Bimodul dan Hasil Kali Tensor}

Ingat bahwa jika $R$ gelanggang komutatif, kita sering cukup mengatakan modul atas $R$
tanpa menyebut apakah modul itu kiri atau kanan.

\begin{definition}
Jika $k$ suatu gelanggang \emph{komutatif}, kita mendefinisikan
\emph{$k$-aljabar} sebagai gelanggang $R$ (yang tidak harus komutatif) beserta
homomorfisme gelanggang $\pi : k \to R$ sedemikian sehingga $\pi(k) \subseteq Z(R)$.
Sebuah \emph{homomorfisme $k$-aljabar} $f : R \to S$ hanyalah homomorfisme
gelanggang sedemikian sehingga $\pi_S = f \circ \pi_R$.
\end{definition}

Secara ekuivalen, $k$-aljabar sekaligus merupakan modul atas $k$ dan gelanggang
yang kedua strukturnya kompatibel.
Perhatikan bahwa pencirian alternatif ini gagal jika $k$ tidak berada
dalam pusat gelanggang perluasan $R$, karena tidak jelas apakah
$R$ seharusnya menjadi modul kiri atau kanan atas $R$.
Jadi, pembatasan pada gelanggang dasar komutatif di pusat
gelanggang perluasan cukup beralasan.

Atas gelanggang nonkomutatif, terdapat beberapa kehalusan pada homomorfisme
dan hasil kali tensor yang dapat diabaikan dalam konteks komutatif.
Di sini kita membahas beberapa kehalusan tersebut.

\begin{definition}
Diberikan gelanggang $R$, \emph{gelanggang lawan} $R^{\mathrm{op}}$ mempunyai
grup abelian dasar yang sama, tetapi perkaliannya berlangsung dalam
urutan terbalik; dengan kata lain, $a \cdot_{\mathrm{op}} b := ba$.
\end{definition}

Perhatikan bahwa gelanggang komutatif secara kanonik isomorfik dengan
gelanggang lawannya melalui peta identitas.
Karena gelanggang dasar $k$ dari $k$-aljabar selalu komutatif,
gelanggang lawan dari suatu aljabar (yang mungkin nonkomutatif)
atas $k$ tetap merupakan $k$-aljabar.

Peta transpos $A \mapsto A^T$ memberikan isomorfisme kanonik
antara gelanggang matriks $\operatorname{M}_n(k)$ dan gelanggang lawannya
$\operatorname{M}_n(k)^{\mathrm{op}}$.
Peta invers $g \to g^{-1}$ memberikan isomorfisme kanonik
antara gelanggang grup $\mathbb{Z}G$ dan gelanggang lawannya
$\mathbb{Z}G^{\mathrm{op}}$.

Sebuah contoh penting bagi tujuan kita menunjukkan bahwa secara umum
kita tidak dapat mengharapkan isomorfisme kanonik:

\begin{example}
Misalkan $V$ ruang vektor berdimensi hingga dengan dual $V^\vee$,
dan perhatikan gelanggang $\operatorname{End}(V)$.
Kita mempunyai isomorfisme kanonik
\[
\psi : \operatorname{End}(V)^{\mathrm{op}} \to \operatorname{End}(V^\vee)
\]
melalui $\psi(f)(g):=g \circ f$ untuk $f \in
\operatorname{End}(V)^{\mathrm{op}}$ dan $g \in V^\vee$.
Kita memperoleh isomorfisme (nonkanonik) antara
$\operatorname{End}(V)^{\mathrm{op}}$ dan $\operatorname{End}(V)$
melalui pilihan isomorfisme $V \cong V^\vee$.
\end{example}

Meskipun contoh-contoh tadi demikian, gelanggang tidak harus isomorfik dengan lawannya.
Kelak kita akan melihat bahwa hal ini bahkan tidak berlaku bagi \emph{gelanggang
divisi}, yang penting bagi teori representasi.
Berikut contoh kecil yang agak dibuat-buat bagi pembaca yang tidak sabar: 

\begin{exercise}
Perhatikan subgelanggang $R \subseteq \operatorname{M}_2(\mathbb{Q})$
yang diberikan oleh
\[
R = \left\{ \begin{pmatrix} a & b\\ 0 & c \end{pmatrix}\ \middle|\
a \in \mathbb{Z}, b \in \mathbb{Q}, c \in \mathbb{Q} \right\}.
\]
Buktikan bahwa $R$ tidak isomorfik dengan $R^{\mathrm{op}}$.
\end{exercise}

Amati bahwa setiap modul kiri atas $R$ merupakan modul kanan atas $R^{\mathrm{op}}$
melalui perkalian skalar $m \cdot_{\mathrm{op}} r := rm$.
Demikian pula, setiap modul kanan atas $R$ merupakan modul kiri atas $R^{\mathrm{op}}$.
Khususnya, seperti sudah kita amati,
pembedaan antara modul kiri dan kanan
atas $R$ tidak berpengaruh tepat ketika $R$ komutatif.

\begin{definition}
Misalkan $R$ dan $S$ gelanggang. Sebuah \emph{$(R,S)$-bimodul} $M$ adalah modul
kiri atas $R$ yang juga merupakan modul kanan atas $S$, sedemikian sehingga $(rm)s=r(ms)$ untuk
semua $r$.
\end{definition}

\begin{exercise}
Tunjukkan bahwa menjadi $(R,S)$-bimodul ekuivalen dengan menjadi
modul kiri atas $R \times S^{\mathrm{op}}$, dan juga dengan menjadi
modul kanan atas $R^{\mathrm{op}} \times S$.
\end{exercise}

Bimodul merupakan objek yang cukup alami:

\begin{example}
Gelanggang $R$ khususnya merupakan $(R,R)$-bimodul.
\end{example}

\begin{example}
Misalkan $V$ dan $W$ modul atas $k$. Himpunan homomorfisme
$\operatorname{Hom}_k(V,W)$ merupakan
$(\operatorname{End}_k(W),\operatorname{End}_k(V))$-bimodul
melalui $fgh := f \circ g \circ h$.
\end{example}

Setiap modul kiri atas $R$ mempunyai struktur kanonik sebagai
$(R,\mathbb{Z})$-bimodul; bahkan sebagai $(R,Z(R))$-bimodul, dengan
$Z(R)$ pusat dari $R$.
Demikian pula, setiap modul kanan atas $R$ mempunyai struktur kanonik sebagai
$(Z(R),R)$-bimodul.
Akibatnya, jika $R$ suatu $k$-aljabar, maka setiap modul kiri atau kanan
atas $R$ secara kanonik sekaligus merupakan modul kiri dan kanan atas $k$.

\begin{definition}
Jika $M$ dan $N$ merupakan modul kiri atas $R$, maka
\[
\operatorname{Hom}_R(M,N) = \operatorname{Hom}_{R-\operatorname{mod}}(M,N)
\]
adalah himpunan homomorfisme modul kiri atas $R$, yaitu $f : M \to N$.
Jika $M$ dan $N$ merupakan modul kanan atas $R$, maka
\[
\operatorname{Hom}_{\operatorname{mod}-R}(M,N)
\]
adalah himpunan homomorfisme modul kanan atas $R$, yaitu $f : M \to N$.
\end{definition}

Misalkan $S$ dan $T$ gelanggang,
$M$ suatu $(R,S)$-bimodul, dan $N$ suatu $(R,T)$-bimodul.
Maka $\operatorname{Hom}_R(M,N)$ mempunyai struktur $(S,T)$-bimodul
melalui $(s\phi t)(m):=\phi(ms)t$ untuk $s \in S, m \in M, t\in T$.
Demikian pula, jika $M$ suatu $(S,R)$-bimodul dan $N$ suatu $(T,R)$-bimodul,
maka $\operatorname{Hom}_{\operatorname{mod}-R}(M,N)$
mempunyai struktur $(T,S)$-bimodul.

Jadi, ketika $R$ komutatif, $\operatorname{Hom}_R(M,N)$
mempunyai struktur modul kanonik atas $R$ berkat struktur
$(R,R)$-bimodul kanonik pada $M$ (dan $N$).
Secara lebih umum, $\operatorname{Hom}_R(M,N)$ hanya mempunyai struktur modul atas $Z(R)$
saja.
Jika $R$ suatu $k$-aljabar, maka $\operatorname{Hom}_R(M,N)$
secara kanonik merupakan modul atas $k$.
Setidaknya, $\operatorname{Hom}_R(M,N)$ selalu merupakan
grup abelian berkat struktur modul kanonik atas $\mathbb{Z}$.

\begin{example}
Misalkan $V$ dan $W$ adalah representasi grup hingga $G$
atas medan $k$.
Maka
\[
\operatorname{Hom}_{kG}(V,W) = \operatorname{Hom}^G_k(V,W)
\]
adalah ruang vektor atas $k$ dari transformasi linear yang $G$-ekuivarian
$f : V \to W$.
\end{example}

\begin{definition}
Misalkan $M$ suatu modul kanan atas $R$, $N$ suatu modul kiri atas $R$,
dan $A$ suatu grup abelian.
Homomorfisme grup $f: M \times N \to A$ disebut \emph{seimbang atas $R$}
jika $f(mr,n)=f(m,rn)$ untuk semua $r \in R$, $m \in M$, dan $n \in N$.
Sebuah \emph{hasil kali tensor} $M \otimes_R N$ adalah grup abelian
beserta homomorfisme grup yang seimbang atas $R$,
$\pi : M \times N \to M \otimes_R N$, sedemikian sehingga
untuk setiap homomorfisme grup yang seimbang atas $R$,
$\phi : M \times N \to A$, terdapat homomorfisme grup tunggal
$\psi :  M \otimes_R N \to A$
sedemikian sehingga $\phi = \psi \circ \pi$.
\end{definition}

Sekali lagi, kita mempunyai:

\begin{proposition}
Hasil kali tensor ada dan tunggal hingga isomorfisme tunggal.
\end{proposition}

Seperti himpunan homomorfisme, hasil kali tensor membawa struktur
tambahan ketika unsur-unsur pembentuknya bimodul.
Jika $M$ suatu $(S,R)$-bimodul dan $N$ suatu $(R,U)$-bimodul,
maka $M \otimes_R N$ merupakan $(S,U)$-bimodul melalui perkalian
$s(m\otimes n)r := (sm) \otimes (nr)$ yang diperluas secara linear.

Banyak konstruksi memakai berbagai struktur bimodul kanonik
secara tersirat.
Sebagai contoh, jika $R$ komutatif, dan $M$ serta $N$ keduanya merupakan
modul \emph{kiri} atas $R$, maka kita dapat memakai struktur
modul \emph{kanan} kanonik atas $R$ pada $M$
untuk memaknai $M \otimes_R N$.
Kemudian $M \otimes_R N$ mempunyai struktur modul atas $R$ dengan memakai struktur
modul kiri atas $R$ pada $M$ atau struktur modul kanan atas $R$ pada $N$
(yang berimpit).
Sebagai contoh lain, jika $R$ suatu $k$-aljabar, maka $M \otimes_R N$
mempunyai struktur modul kanonik atas $k$.

Sekarang kita siap menyatakan suatu hasil penting:

\begin{theorem}[Adjoin Tensor--Hom]
Misalkan $R$, $S$, $U$, $V$ gelanggang.
Misalkan $M$ suatu $(R,S)$-bimodul, $N$ suatu $(S,U)$-bimodul,
dan $P$ suatu $(R,V)$-bimodul.
Maka terdapat isomorfisme alami
\[
\operatorname{Hom}_R(M \otimes_S N,P)
\cong \operatorname{Hom}_S(N,\operatorname{Hom}_R(M,P))
\]
antara $(U,V)$-bimodul.
\end{theorem}

\begin{exercise}
Himpunan homomorfisme dalam pernyataan Adjoin Tensor--Hom
memakai struktur modul \emph{kiri} pada $M,N,P$, dan seterusnya.
Tentukan pernyataan serupa ketika himpunan homomorfisme tersebut
merujuk pada struktur modul \emph{kanan}.
\end{exercise}

Jarang sekali semua struktur bimodul diperlukan sekaligus!
Kasus khusus yang terpenting terjadi ketika $M$ suatu $(R,S)$-bimodul,
$N$ suatu modul kiri atas $S$, dan $P$ suatu modul kiri atas $R$;
isomorfisme yang dihasilkan ketika itu hanyalah isomorfisme grup abelian.

Salah satu penerapan utama Adjoin Tensor--Hom adalah
memahami restriksi dan perluasan skalar.

\begin{definition}
Misalkan $f : R \to S$ suatu homomorfisme gelanggang.
Jika $N$ suatu modul kiri atas $S$, maka \emph{restriksi skalar
dari $N$} adalah struktur modul kiri atas $R$ pada $N$ yang diberikan oleh
$r \cdot_R n := f(r)n$ untuk $r \in R$ dan $n \in N$.
Kita sering menyatakan restriksi skalar dengan ${}_RN$
atau ${}_fN$ untuk menekankan pembedaan dengan $N$ semula
yang dipandang sebagai modul kiri atas $S$.
\end{definition} 

Ketika $f$ merupakan inklusi subgelanggang, restriksi
skalar dapat dipandang sebagai “melupakan” sebagian struktur $N$.
Sebagai contoh, ruang vektor kompleks $\mathbb{C}^n$ menjadi ruang
vektor real $\mathbb{R}^{2n}$ melalui restriksi oleh inklusi
$\mathbb{R} \to \mathbb{C}$.

\begin{example}
Jika $H$ subgrup dari grup $G$, kita mempunyai inklusi
$kH \to kG$ antara aljabar grup.
Misalkan $V$ suatu representasi grup hingga $G$ atas medan $k$.
Maka $V$ mempunyai struktur modul kiri atas $kG$.
Restriksi skalar ${}_{kH}V$ adalah modul kiri atas $kH$
yang bersesuaian dengan $\operatorname{Res}_H^G V$. 
\end{example}

Beralih ke arah sebaliknya sedikit lebih rumit.

\begin{definition}
Misalkan $f : R \to S$ suatu homomorfisme gelanggang dan $M$ suatu modul atas $R$.
Misalkan $S_R$ adalah struktur $(S,R)$-bimodul yang bersesuaian pada $S$,
dan misalkan ${}_RS$ adalah struktur $(R,S)$-bimodul yang bersesuaian pada $S$.
\emph{Perluasan skalar dari $M$}
adalah modul kiri atas $S$, yaitu $S_R \otimes_R M$.
\emph{Koperluasan skalar dari $M$}
adalah modul kiri atas $S$, yaitu $\operatorname{Hom}_R({}_RS,M)$.
\end{definition}

Perhatikan bahwa notasi $S_R$ dan ${}_RS$ biasanya dianggap
terlalu terperinci.
Kita biasanya cukup menulis $S \otimes_R M$ untuk perluasan skalar
dan $\operatorname{Hom}_R(S,M)$ untuk koperluasan skalar.
Notasi ini biasanya tak taksa, walaupun mungkin membingungkan bagi pemula.

Adjoin Tensor--Hom menunjukkan bahwa perluasan skalar
dan restriksi skalar saling adjoin.
Memang, dengan memandang $S$ sebagai $(S,R)$-bimodul,
kita mempunyai $\operatorname{Hom}_S(S,M) \cong {}_RM$.
Jadi, adjoin tersebut menjadi
\[
\operatorname{Hom}_R(S \otimes_R N,M)
\cong \operatorname{Hom}_R(N,{}_RM)
\]
untuk modul kiri atas $R$, yaitu $N$, dan modul kiri atas $S$, yaitu $M$.

Sekali lagi, hal-hal ini langsung berkaitan dengan teori representasi:

\begin{example}
Jika $H$ subgrup dari grup $G$, kita mempunyai inklusi
$kH \to kG$ antara aljabar grup.
Misalkan $W$ suatu representasi grup hingga $H$ atas medan $k$.
Maka $W$ mempunyai struktur modul kiri atas $kH$.
Perluasan skalar $kG \otimes_{kH} W$
adalah modul kiri atas $kG$ yang bersesuaian dengan $\operatorname{Ind}_H^G W$.
\end{example} 

Untuk melihat mengapa perluasan skalar bersesuaian dengan induksi, kita hanya perlu
mengingat bahwa kita pada dasarnya \emph{mendefinisikan} induksi sebagai representasi
linear yang memenuhi (suatu versi) resiprositas Frobenius.
Adjoin Tensor--Hom memberikan isomorfisme ruang vektor atas $k$
\[
\operatorname{Hom}_{kG}(kG \otimes_{kH} W,V)
\cong \operatorname{Hom}_{kH}(W,{}_{kH}V))
\]
yang tepat merupakan isomorfisme
\[
\operatorname{Hom}_k^G(\operatorname{Ind}_H^G W,V)
\cong \operatorname{Hom}_k^H(W,\operatorname{Res}_H^G V)
\]
dari Resiprositas Frobenius!

Kita serahkan sebagai latihan pemeriksaan bahwa koperluasan skalar
bersesuaian dengan koinduksi. Memang, pembedaan itu tidak penting
(setidaknya dalam konteks kita) berdasarkan pernyataan berikut:

\begin{exercise}
Misalkan $H$ subgrup dari grup hingga $G$, dan $W$ suatu modul atas $kH$.
Buktikan bahwa
\[
kG \otimes_{kH} W \cong \operatorname{Hom}_{kH}(kG,W)
\]
sebagai modul kiri atas $kG$.
\end{exercise}

\section{Dekomposisi Wedderburn}

Sepanjang bagian ini, $k$ adalah medan.
Tujuan utama bagian ini ialah membuktikan hasil berikut:

\begin{theorem}[Teorema Wedderburn] \label{thm:Wedderburn}
Sebuah $k$-aljabar berdimensi hingga $A$ bersifat \emph{semisederhana}
jika dan hanya jika
\[
A \cong \bigoplus_{i=1}^r \operatorname{M}_{n_i}(D_i)
\]
dengan $D_1, \ldots, D_r$ merupakan $k$-aljabar berdimensi hingga
dan $n_1,\ldots, n_r$ bilangan bulat positif.
\end{theorem}

\begin{definition}
Misalkan $A$ suatu $k$-aljabar dan $M$ suatu modul kiri atas $A$.
Sebuah \emph{deret komposisi} atau \emph{filtrasi sederhana hingga}
adalah barisan menurun submodul
\[
M=M_0 \supsetneq M_1 \supsetneq \cdots \supsetneq M_r = 0
\]
sedemikian sehingga setiap hasil bagi $M_i/M_{i+1}$ sederhana.
Bilangan bulat $r$ adalah \emph{panjang} filtrasi tersebut.
Modul $M$ disebut \emph{berpanjang hingga} jika memiliki filtrasi sederhana
hingga atau $M=0$.
\end{definition}

\begin{theorem}[Jordan-H\"older]
Setiap dua deret komposisi dari modul kiri $M$ bersifat \emph{ekuivalen}:
kelas-kelas isomorfisme hasil bagi $M_i/M_{i+1}$
tunggal hingga pengurutan ulang.
\end{theorem}

Berdasarkan teorema Jordan-H\"older, pengertian berikut
terdefinisi dengan baik:

\begin{definition}
\emph{Panjang} modul $M$ berpanjang hingga adalah panjang dari sebarang
deret komposisinya (atau $0$ jika $M=0$).
\end{definition}

Ingat bahwa gelanggang disebut Artin (berturut-turut, Noether) jika memenuhi
syarat rantai menurun (berturut-turut, syarat rantai menaik) pada ideal.
Demikian pula, modul disebut Artin (berturut-turut, Noether) jika memenuhi
syarat rantai menurun (berturut-turut, syarat rantai menaik) pada
submodul. Khususnya, ruang vektor atas medan bersifat Artin
sekaligus Noether, sehingga kita langsung memperoleh pernyataan berikut.

\begin{proposition}
Misalkan $A$ suatu $k$-aljabar berdimensi hingga,
dan $M$ suatu modul kiri atas $A$ yang terbangkitkan hingga.
Maka $A$ dan $M$ keduanya bersifat Noether dan Artin.
Khususnya, kita melihat bahwa $M$ berpanjang hingga.
\end{proposition}

Jika $M$ suatu modul kiri semisederhana, maka terdapat dekomposisi
\begin{equation} \label{eq:isoDecomp}
M = M_1 \oplus \cdots \oplus M_r,
\end{equation}
dengan setiap $M_i$ submodul yang isomorfik dengan $S_i^{\oplus a_i}$,
$S_i$ sederhana, $a_i$ bilangan bulat positif,
dan $S_i \not\cong S_j$ ketika $i \ne j$.
Perhatikan bahwa setiap submodul $M_i$ bersifat kanonik, tetapi
isomorfisme $M_i \cong S_i^{\oplus a_i}$ tidak.
Submodul-submodul $M_1, \ldots, M_r$ disebut \emph{komponen
isotipik} dari $M$, dan \eqref{eq:isoDecomp} disebut \emph{dekomposisi
isotipik} dari $M$.

\begin{lemma} \label{lem:ssModule}
Misalkan $A$ suatu $k$-aljabar berdimensi hingga dan $M$ suatu modul kiri
atas $A$ yang terbangkitkan hingga.
Pernyataan-pernyataan berikut ekuivalen:
\begin{enumerate}
\item[(a)] setiap submodul $M$ merupakan suku langsung,
\item[(b)] $M$ semisederhana, dan
\item[(c)] $M$ merupakan jumlah submodul-submodul sederhana (belum tentu langsung).
\end{enumerate}
\end{lemma}

\begin{proof}
Jika $N_1 \subset N_2 \subset M$ suatu rantai submodul dan $N_1$ merupakan
suku langsung dari $M$, maka $N_1$ merupakan suku langsung dari $N_2$.
Jadi, (a) $\implies$ (b) diperoleh melalui induksi pada dimensi $M$,
dengan modul sederhana sebagai kasus dasar. Implikasi (b) $\implies$
(c) langsung. Tinggal menunjukkan (c) $\implies$ (a).

Misalkan $N$ suatu submodul dari $M$.
Misalkan $N'$ submodul maksimal dari $M$ sedemikian sehingga $N \cap N' = 0$
(yang ada karena $M$ berdimensi hingga).
Jika $N + N' = M$, pembuktian selesai; jadi, andaikan $N + N' \subsetneq M$.
Maka terdapat submodul sederhana $S$ dari $M$ yang tidak termuat dalam $N+N'$;
bahkan, $S \cap (N+N')=0$ karena submodul itu sederhana.

Khususnya, $N' \subsetneq N'+S$.
Perhatikan $m \in (N'+S) \cap N$. Kita mempunyai $m=n+s$, dengan $n \in N'$
dan $s \in S$. Unsur $s=m-n$ termuat dalam $S$ sekaligus $N+N'$, sehingga
nol. Jadi, $m \in N' \cap N$ juga nol.
Kita menyimpulkan bahwa $(N'+S) \cap N$ nol.
Modul $N'+S$ bertentangan dengan asumsi kemaksimalan
$N'$.
\end{proof}

\begin{lemma} \label{lem:subquot}
Submodul dan modul hasil bagi dari modul semisederhana bersifat semisederhana.
\end{lemma}

\begin{proof}
Submodul bersifat semisederhana melalui argumen yang serupa dengan
(1) $\implies$ (2) dari Lema~\ref{lem:ssModule}.
Modul hasil bagi bersifat semisederhana berdasarkan lema Schur dan
syarat (3) dari Lema~\ref{lem:ssModule}.
\end{proof}

Ingat bahwa aljabar $A$ semisederhana jika dan hanya jika merupakan hasil kali
langsung gelanggang-gelanggang sederhana. Kita akan melihat bahwa ini ekuivalen dengan $A$
semisederhana sebagai modul kiri atas dirinya sendiri. Sementara itu, kita membuktikan
rumusan ekuivalen ketiga:

\begin{lemma}
Misalkan $A$ suatu $k$-aljabar berdimensi hingga.
Setiap modul kiri atas $A$ yang terbangkitkan hingga semisederhana jika dan hanya jika
${}_AA$ semisederhana sebagai modul kiri atas $A$ itu sendiri.
\end{lemma}

\begin{proof}
Andaikan ${}_AA$ semisederhana sebagai modul kiri atas $A$.
Setiap modul atas $A$, yaitu $M$, dapat ditulis sebagai hasil bagi modul bebas atas $A$,
$A^n$, untuk suatu bilangan bulat positif $n$.
Jadi, setiap submodul $M$ semisederhana berdasarkan Lema~\ref{lem:subquot}.
Kebalikannya berlaku terlebih lagi.
\end{proof}

Sebagai akibat lema sebelumnya, bersama dengan
Jordan-H\"older, kita menyimpulkan bahwa aljabar semisederhana hanya mempunyai
berhingga banyak kelas isomorfisme modul sederhana!

\begin{lemma}
Misalkan $A$ suatu $k$-aljabar berdimensi hingga.
Jika $A$ aljabar sederhana, maka ${}_AA$ semisederhana.
\end{lemma}

\begin{proof}
Misalkan $M$ jumlah semua submodul sederhana dalam ${}_AA$.
Jika $S$ suatu submodul sederhana dalam ${}_AA$ dan $a \in A$, maka ideal
kiri $Sa$ adalah $0$ atau submodul sederhana dari ${}_AA$.
Jadi, $Ma \subseteq M$ untuk setiap $a \in A$.
Karena itu, $M$ ideal kanan dari $A$, bahkan ideal dua sisi.
Karena $M \ne 0$ dan $A$ sederhana sebagai gelanggang, berlaku $M=A$.
Jadi, ${}_AA$ merupakan jumlah submodul-submodul sederhana dan karena itu semisederhana.
\end{proof}

\begin{corollary} \label{cor:semisimple_def}
Misalkan $A$ suatu $k$-aljabar berdimensi hingga.
Jika $A$ aljabar semisederhana, maka ${}_AA$ semisederhana.
\end{corollary}

\begin{proof}
Misalkan $A = A_1 \oplus \cdots \oplus A_n$, dengan setiap $A_i$ suatu
$k$-aljabar sederhana. Amati bahwa struktur modul kiri atas $A$ pada ${}_AA_i$
memfaktor melalui struktur modul atas $A_i$.
Setiap $A_i$, khususnya, merupakan submodul dari ${}_AA$.
Karena semuanya semisederhana, jumlahnya juga semisederhana.
\end{proof}

(Kebalikan korolari ini akan diperoleh sebagai akibat
teorema Wedderburn.)

Bukti pernyataan berikut diserahkan sebagai latihan:

\begin{lemma} \label{lem:fancyMatrix}
Misalkan $A$ suatu gelanggang.
Misalkan $M= M_1 \oplus \cdots \oplus M_m$ dan
$N=N_1 \oplus \cdots \oplus N_n$ merupakan jumlah langsung modul-modul atas $A$.
Maka
\[ \operatorname{Hom}_A\left(\bigoplus_iM_i,\ \bigoplus_jN_i\right)
\cong \bigoplus_{i,j}\operatorname{Hom}_A(M_i, N_j) \ . \]
\end{lemma}

\begin{lemma}
Misalkan $A$ suatu $k$-aljabar berdimensi hingga, dan
misalkan $M$ suatu modul kiri semisederhana atas $A$.
Andaikan
\[
M = M_1 \oplus \cdots \oplus M_r
\]
adalah dekomposisi isotipik dengan $M_i=S_i^{\oplus n_i}$
untuk suatu modul sederhana atas $A$, yaitu $S_i$.
Terdapat isomorfisme kanonik $k$-aljabar
\[
\operatorname{End}_A(M) \cong \operatorname{M}_{n_1}(\operatorname{End}_A(S_1))
\oplus \cdots \oplus
\operatorname{M}_{n_r}(\operatorname{End}_A(S_r)) .
\]
\end{lemma}

\begin{proof}
Berdasarkan Lema~\ref{lem:fancyMatrix},
kita mempunyai $\operatorname{End}_A(N) = \operatorname{Hom}_A(N,N)$ dan
$\operatorname{End}_A(M^{\oplus n}) \cong M_n(\operatorname{End}_A(N))$
untuk setiap modul kiri atas $A$, yaitu $N$, dan setiap bilangan bulat positif $n$.
Jadi, cukup
ditunjukkan bahwa $\operatorname{Hom}_A(M_i,M_j)=0$ setiap kali $i \ne j$.
Perhatikan bahwa $M_i=S_i^{\oplus n_i}$ dan $M_j=S_j^{\oplus n_j}$
untuk bilangan bulat $n_i, n_j$ dan
modul sederhana tak isomorfik $S_i$ dan $S_j$.
Dengan kembali memakai lema sebelumnya, $\operatorname{Hom}_A(M_i,M_j)$ merupakan jumlah langsung
modul-modul berbentuk $\operatorname{Hom}_A(S_i,S_j)$. Modul-modul ini bernilai $0$ berdasarkan lema Schur.
\end{proof}

\begin{lemma}
Misalkan $A$ suatu $k$-aljabar.
Terdapat isomorfisme kanonik $A^{\mathrm{op}} \cong \operatorname{End}_A({}_AA)$ antara
$k$-aljabar.
\end{lemma}

\begin{proof}
Diberikan unsur $a \in A^{\mathrm{op}}$,
terdapat endomorfisme tunggal $\phi_a \in A^{\mathrm{op}}$
sedemikian sehingga $\phi_a(1)=a$.
Ini memberikan korespondensi yang komutatif dengan penjumlahan dan
perkalian skalar.
Jadi, terdapat isomorfisme antara ruang-ruang vektor yang mendasarinya.
Perhitungan
\[ (\phi_a \circ \phi_b)(1)=\phi_a(b)=b\phi_a(1)=ba \]
menunjukkan bahwa perkaliannya berimpit.
Jadi, kedua aljabar tersebut isomorfik.
\end{proof}

Satu pengamatan lagi sebelum teorema utama:

\begin{lemma}
Untuk $k$-aljabar $A$, terdapat
isomorfisme $k$-aljabar $\operatorname{M}_n(A)^{\mathrm{op}}
\cong \operatorname{M}_n(A^{\mathrm{op}})$.
\end{lemma}

\begin{proof}
Petakan suatu matriks ke transposnya. Ini jelas merupakan
isomorfisme ruang vektor. Perhitungan singkat menunjukkan bahwa perkaliannya
juga kompatibel. 
\end{proof}

Sekarang kita membuktikan teorema Wedderburn.

\begin{proof}[Bukti Teorema~\ref{thm:Wedderburn}]
Andaikan $A$ semisederhana sebagai $k$-aljabar.
Maka ${}_AA$ semisederhana sebagai modul kiri atas $A$
berdasarkan Korolari~\ref{cor:semisimple_def}.

Karena ${}_AA$ semisederhana, ia mempunyai dekomposisi isotipik
\[
{}_AA = S_1^{\oplus n_1} \oplus \cdots \oplus S_r^{\oplus n_r}
\]
dengan $S_i$ modul-modul sederhana yang dua-dua tidak isomorfik.
Misalkan $D_i=\operatorname{End}_A(S_i)$,
yang merupakan aljabar divisi berdasarkan Lema Schur.

Dengan demikian,
\begin{align*}
A^{\operatorname{op}} &\cong \operatorname{End}_A(A)\\
&\cong \operatorname{End}_A(S_1^{\oplus n_1}) \oplus \cdots
\oplus \operatorname{End}_A(S_r^{\oplus n_r})\\
&\cong M_{n_1}(D_1) \oplus \cdots \oplus M_{n_r}(D_r) \ ,
\intertext{dan dengan menerapkan lema terakhir:}
A &\cong M_{n_1}(D_1^{\operatorname{op}}) \oplus \cdots \oplus
M_{n_r}(D_r^{\operatorname{op}}) \ .
\end{align*}

Kebalikannya kita serahkan sebagai latihan.
\end{proof}

Perhatikan bahwa korolari dari uraian di atas (yang salah satu arahnya dipakai
dalam bukti) adalah pencirian berikut bagi kesemisederhanaan suatu
aljabar.

\begin{corollary}
Misalkan $A$ suatu $k$-aljabar berdimensi hingga.
Pernyataan-pernyataan berikut ekuivalen:
\begin{itemize}
\item $A$ semisederhana sebagai gelanggang (suatu jumlah langsung gelanggang sederhana).
\item $A$ semisederhana sebagai modul kiri atas $A$ (suatu jumlah langsung modul
sederhana).
\item Setiap modul kiri atas $A$ yang terbangkitkan hingga semisederhana.
\end{itemize}
\end{corollary}

\section{Aljabar Sederhana Sentral}

Seperti dalam bagian sebelumnya, $k$ adalah medan.
Kita menyatakan penutupan aljabar $k$ dengan $\overline{k}$.
Beberapa hasil di sini diuraikan lebih lanjut dalam teks tentang aljabar sederhana
sentral dan grup Brauer (misalnya, \cite{GS} atau \cite{Saltman}). 

Karena aljabar semisederhana merupakan jumlah langsung aljabar sederhana,
sekarang kita tertarik mempelajari aljabar sederhana.
Invarian terpenting suatu aljabar sederhana adalah pusatnya,
yang selalu merupakan medan; jika tidak, terdapat ideal dua sisi
nontrivial. Hal ini memotivasi definisi berikut.

\begin{definition}
Sebuah \emph{$k$-aljabar sederhana sentral} adalah $k$-aljabar berdimensi hingga
$A$ dengan pusat $Z(A)=k$.
Sebuah \emph{$k$-aljabar divisi sentral} adalah $k$-aljabar berdimensi hingga
$A$ dengan pusat $Z(A)=k$.
\end{definition}

Ingat pernyataan berikut, yang kita buktikan ketika menegakkan Lema Schur
dalam kasus kompleks:

\begin{proposition} \label{prop:division_ac}
Jika $k$ tertutup secara aljabar, maka setiap $k$-aljabar divisi sentral
isomorfik dengan $k$.
\end{proposition}

Berdasarkan teorema Wedderburn, kita mempunyai deskripsi yang baik tentang aljabar
sederhana sentral dan hubungannya dengan jenis-jenis
aljabar lain yang sudah kita lihat. Jika $A$ suatu $k$-aljabar berdimensi hingga,
maka:

\begin{enumerate}
\item Jika $A$ semisederhana, maka aljabar itu merupakan hasil kali $k$-aljabar sederhana.
\item Jika $A$ sederhana, maka aljabar itu merupakan $K$-aljabar sederhana sentral, dengan
$K$ suatu perluasan medan hingga dari $k$.
\item Jika $A$ suatu $k$-aljabar sederhana sentral, maka $A \cong
\operatorname{M_n}(D)$, dengan $D$ suatu $k$-aljabar divisi sentral.
\end{enumerate}

Aljabar sederhana sentral dan aljabar divisi berhubungan erat.
Dalam suatu arti, jenis pertama dapat dipahami sepenuhnya melalui jenis kedua:

\begin{proposition}
Jika $A$ suatu $k$-aljabar sederhana sentral, maka terdapat aljabar divisi sentral tunggal
atas $k$, yaitu $D$ (hingga isomorfisme), dan bilangan bulat positif $n$
sedemikian sehingga $A \cong \operatorname{M}_n(D)$.
\end{proposition}

\begin{proof}
Andaikan $A \cong M_n(D) \cong M_m(E)$ untuk $k$-aljabar divisi sentral
$D$ dan $E$.
Misalkan $L$ suatu ideal kiri minimal dari $A$.
Maka $D^{\oplus n} \cong L \cong E^{\oplus m}$ karena semua aljabar
tersebut sederhana dan hanya mempunyai satu kelas isomorfisme modul sederhana.
Kita dapat memandang $D^{\oplus n}$ sebagai himpunan vektor kolom dengan
entri-entri dalam $D$.
Amati bahwa endomorfisme $D^{\oplus n} \to D^{\oplus n}$
sebagai modul kiri atas $A$ sepenuhnya ditentukan oleh citra
unsur $(1,0,\ldots,0)$. Dapat diperiksa bahwa satu-satunya
citra yang diperbolehkan berbentuk $(a,0,\ldots,0)$ untuk $a \in A$.
Jadi, $\operatorname{End}_A(D^{\oplus n}) \cong D$.
Dengan demikian,
\[
D \cong \operatorname{End}_A(D^{\oplus n})
\cong L
\cong \operatorname{End}_A(E^{\oplus m}) \cong E
\]
dan $m=n$ diperoleh dengan menghitung dimensi.
\end{proof}

Namun, aljabar sederhana sentral muncul secara alami dalam konstruksi-konstruksi
yang agak canggung jika dibatasi pada aljabar divisi saja.

\begin{lemma} \label{lem:center_of_tensor}
Jika $A$ dan $B$ merupakan $k$-aljabar, maka $Z(A \otimes_k B) = Z(A) \otimes_k
Z(B)$.
\end{lemma}

\begin{proof}
Jelas bahwa $Z(A) \otimes_k Z(B) \subseteq Z(A \otimes_k B)$.
Sekarang, andaikan $z = \sum_i a_i \otimes b_i$ suatu unsur
$Z(A \otimes_k B)$. Kita dapat mengasumsikan $a_i$ bebas linear.
Ambil $x \in B$ dan amati bahwa
\[
0 = z(1 \otimes x) - (1 \otimes x)z =
\sum_i a_i \otimes (b_ix-xb_i),
\]
yang mengakibatkan setiap $b_ix-xb_i=0$.
Jadi, setiap $b_i \in Z(B)$.
Dengan demikian, $z \in A \otimes_k Z(B)$.
Sekarang, kita dapat mengasumsikan (setelah mungkin menulis ulang $z$) bahwa $b_i$
bebas linear dalam $Z(B)$
dan serupa dengan itu menyimpulkan bahwa setiap $a_i$ berada dalam $Z(A)$.
\end{proof}

\begin{lemma} \label{lem:csa_simple_simple}
Jika $A$ suatu $k$-aljabar sederhana sentral dan
$B$ suatu $k$-aljabar sederhana, maka $A \otimes_k B$ merupakan
$k$-aljabar sederhana.
\end{lemma}

\begin{proof}
Misalkan $I$ suatu ideal dua sisi nontrivial dalam $A \otimes_k B$.
Misalkan $r$ bilangan bulat positif minimal sedemikian sehingga
terdapat unsur $x \in I$ yang dapat ditulis
\[
x = \sum_{i=1}^r a_i \otimes b_i.
\]
Karena $I$ berada dalam $A$ yang sederhana, $Aa_1A=A$. Jadi, kita mempunyai $y,z \in A \otimes 1$
sedemikian sehingga
\[
x'=yxz = 1 \otimes b_1 + \sum_{i=2}^r a_i' \otimes b_i
\]
untuk $a_i'=ya_iz$, dengan $x' \in I$ pula.
Untuk semua $w \in A$, kita melihat bahwa
\[
(w \otimes 1)x' - x'(w \otimes 1) = -\sum_{i=2}^r (wa_i'-a_i'w) \otimes
b_i = 0
\]
berdasarkan keminimalan $r$. Jadi, $x' \in Z(A) \otimes_k B = k \otimes_k B$.
Dengan demikian, $x'=1 \otimes b \in I$ untuk suatu $b \in B$.
Namun, $B$ sederhana, sehingga $BbB=B$. Maka $1 \otimes 1 \in I$.
Kita menyimpulkan bahwa $A \otimes B = I$, seperti diinginkan.
\end{proof}

Perhatikan bahwa hipotesis bahwa $A$ sentral atas $k$ sangat penting
dalam Lema~\ref{lem:csa_simple_simple} --- perhatikan $\mathbb{C}
\otimes_{\mathbb{R}} \mathbb{C}$.

\begin{lemma}
Jika $A$ suatu $k$-aljabar sederhana sentral berdimensi hingga,
maka terdapat isomorfisme $k$-aljabar
$A \otimes_k A^{\mathrm{op}} \cong M_n(k)$
untuk $n = \dim_k(A)$.
\end{lemma}

\begin{proof}
Dengan memperluas secara linear, terdapat homomorfisme $k$-aljabar
$\Psi : A \otimes_k A^{\mathrm{op}} \to \operatorname{End}_k(A)$
sedemikian sehingga $\Psi(x \otimes y)(z)=xzy$ untuk $x,z \in A$ dan
$y \in A^{\mathrm{op}}$. Peta tersebut injektif karena
$A \otimes A^{\mathrm{op}}$ sederhana.
Kesurjektifan diperoleh dengan menghitung dimensi.
\end{proof}

(Sebenarnya, kebalikan lema sebelumnya juga benar, tetapi kita tidak
memerlukannya.)

\begin{definition}
Dua $k$-aljabar sederhana sentral $A$ dan $B$ disebut
\emph{ekuivalen Brauer} jika $A \otimes_k M_n(k) \cong B \otimes_k
M_m(k)$ untuk suatu bilangan bulat positif $m,n$.
\emph{Grup Brauer}, yang dinyatakan dengan $\operatorname{Br}(k)$,
adalah himpunan kelas ekuivalensi Brauer
dari $k$-aljabar sederhana sentral.
\end{definition}

Dari lema-lema di atas, kita memperoleh:

\begin{proposition}
Grup Brauer $\operatorname{Br}(k)$ merupakan grup terhadap
hasil kali tensor wakil-wakilnya,
dengan identitas berupa kelas aljabar matriks $M_n(k)$
dan invers diberikan oleh $[A] \mapsto [A^{\mathrm{op}}]$.
\end{proposition}

Secara ekuivalen, $A$ dan $B$ ekuivalen Brauer jika $A$ dan $B$ keduanya
merupakan gelanggang matriks atas $k$-aljabar divisi sentral yang sama.
Jadi, $\operatorname{Br}(k)$ juga dapat dipandang sebagai himpunan kelas
isomorfisme $k$-aljabar divisi sentral; namun, struktur grupnya
lebih sukar dilihat dengan definisi ini.

Alasan lain kita mempelajari aljabar sederhana sentral, bukan hanya
aljabar divisi yang mendasarinya, ialah bahwa sifat menjadi aljabar sederhana
sentral invarian terhadap perluasan medan.

\begin{proposition}
Misalkan $A$ suatu $k$-aljabar berdimensi hingga.
Pernyataan-pernyataan berikut ekuivalen:
\begin{enumerate}
\item[(a)] $A$ merupakan $k$-aljabar sederhana sentral.
\item[(b)] $A \otimes_k K$ merupakan $K$-aljabar sederhana sentral untuk suatu perluasan
medan $K/k$.
\item[(c)] $A \cong \operatorname{M}_n(D)$ untuk suatu $k$-aljabar divisi sentral
$D$ dan suatu $n$.
\item[(d)] $A \otimes_k \overline{k} \cong \operatorname{M}_n(\overline{k})$
untuk suatu $n$.
\item[(e)] $A \otimes_k K \cong \operatorname{M}_n(K)$
untuk suatu perluasan medan hingga $K/k$ dan suatu $n$.
\item[(f)] $A \otimes_k K \cong \operatorname{M}_n(K)$
untuk suatu perluasan medan Galois hingga $K/k$ dan suatu $n$.
\end{enumerate}
\end{proposition}

\begin{proof}
Ekuivalensi (e) dan (f) pada dasarnya bergantung pada pembuktian bahwa kita dapat
memilih perluasan medan \emph{separabel}.
Karena kita terutama bekerja dalam
karakteristik $0$, kita tidak akan membahas (f) di sini;
lihat \cite[Proposition 2.2.5]{GS} untuk buktinya.

Kita mulai dengan menegakkan (a) $\implies$ (b).
Berdasarkan Lema~\ref{lem:csa_simple_simple}, kita menyimpulkan bahwa $A \otimes_k K$
sederhana karena $K$ sederhana. Pusat $A \otimes_k K$
adalah $k \otimes_k K=K$ berdasarkan Lema~\ref{lem:center_of_tensor}.
Jadi, $A \otimes_k K$ merupakan $K$-aljabar sederhana sentral, seperti diinginkan.

Sekarang perhatikan (b) $\implies$ (a).
Jika $J$ ideal dua sisi dari $A$, maka $J \otimes_k K$ merupakan
ideal dua sisi dari $A \otimes_k K$; jadi, $A$ sederhana.
Jika $z$ berada dalam pusat $Z(A)$, maka $z \otimes 1$ berada dalam
pusat $Z(A \otimes K)$. Jadi, $z \in K \cap k = k$.

Implikasi (a) $\implies$ (c) sudah kita lihat diperoleh dari
teori Wedderburn.
Implikasi (c) $\implies$ (d) diperoleh dari
Proposisi~\ref{prop:division_ac},
dengan memakai fakta bahwa kita sudah menunjukkan (a) $\implies$ (b).
Implikasi (d) $\implies$ (e) diperoleh dari prinsip Lefschetz:
koefisien-koefisien dalam isomorfisme eksplisit
$A \otimes_k \overline{k} \cong \operatorname{M}_n(\overline{k})$
harus terdefinisi atas suatu perluasan hingga dari $k$.
Implikasi (e) $\implies$ (b) langsung.
\end{proof}

Berdasarkan proposisi sebelumnya, kita dapat membahas beberapa
konsep baku yang terkait dengan aljabar sederhana sentral.

\begin{definition}
Misalkan $A$ suatu $k$-aljabar sederhana sentral.
Aljabar $A$ disebut \emph{terurai} jika $A \cong \operatorname{M}_n(k)$
untuk suatu bilangan bulat $n$. Sebuah \emph{medan pengurai} dari $A$ adalah
perluasan medan $K/k$ sedemikian sehingga $A \otimes_k K \cong \operatorname{M}_n(K)$.
\end{definition}

Perhatikan bahwa jika $V$ suatu ruang vektor atas $k$,
maka $\dim_k(V) = \dim_K(V \otimes_k K)$
untuk semua perluasan medan $K/k$.
Jadi, dimensi $\dim_k(A)$ dari $k$-aljabar sederhana sentral $A$
selalu merupakan kuadrat, karena
$A \otimes_k K \cong \operatorname{M}_n(K)$ untuk suatu perluasan medan
$K/k$.

\begin{definition}
Misalkan $A$ suatu $k$-aljabar sederhana sentral.
\emph{Derajat} dari $A$ adalah bilangan bulat positif tunggal
$\deg(A)$ sedemikian sehingga
$\dim_k(A)=\deg(A)^2$.
\emph{Indeks (Schur)} dari $A$ diberikan oleh $\operatorname{ind}(A)=\deg(D)$,
dengan $D$ suatu $k$-aljabar divisi sentral sedemikian sehingga
$A\cong \operatorname{M}_n(D)$.
\end{definition}

\begin{proposition}
Jika $A$ suatu $k$-aljabar sederhana sentral dan
$M$ suatu modul kiri sederhana atas $A$, maka
$\dim_k(M) =\deg(A) \operatorname{ind}(A)$.
\end{proposition}

\begin{proof}
Perhatikan bahwa $A \cong \operatorname{M}_r(D)$ untuk suatu $k$-aljabar divisi
sentral $D$.
Karena $M$ sederhana, kita mempunyai $M \cong D^{\oplus r}$.
Jadi, $\dim_k(M) = r \deg(D)^2=\deg(A) \operatorname{ind}(A)$.
\end{proof}

Misalkan $A$ suatu $k$-aljabar sederhana sentral dengan medan pengurai $K$,
dan misalkan $\Psi : A \hookrightarrow M_n(K)$ adalah pembenaman
yang diperoleh melalui komposisi dengan isomorfisme $A \otimes_k K \cong
M_n(K)$.
Diberikan unsur $a \in A$, kita mendefinisikan
\emph{polinomial karakteristik tereduksi}
\[
\operatorname{Prd}_{A,a}(t) := \chi_{\Psi(a)}(t)
\]
dengan $\chi_{\Psi(a)}(t)$ polinomial karakteristik
dari $\Psi(a)$ dalam $M_n(K)$.

\begin{proposition}
Polinomial karakteristik tereduksi
$\operatorname{Prd}_{A,a}(t)$ mempunyai koefisien dalam $k$ dan tidak bergantung pada
pilihan medan pengurai $K$ maupun pilihan isomorfisme
$A \otimes_k K \cong M_n(K)$.
\end{proposition}

\begin{proof}
Kita menghilangkan bukti karena hasil ini tidak mendasar bagi alur pembahasan kita.
\end{proof}

Dengan menulis
\[
\operatorname{Prd}_{A,a}(t) = t^n - e_1t^{n-1} + \cdots + (-1)^n e_n
\]
sebagai polinomial dalam $k[t]$, kita mendefinisikan
\emph{teras tereduksi} $\operatorname{Trd}_A(a)=e_1$
dan \emph{norma tereduksi} $\operatorname{Nrd}_A(a)=e_n$.

\begin{proposition}
Misalkan $A$ suatu $k$-aljabar sederhana sentral.
Misalkan $\chi_{A/k,a}(t)$, $\operatorname{Tr}_{A/k}(a)$, dan
$\operatorname{N}_{A/k}(a)$
berturut-turut merupakan polinomial karakteristik, teras, dan norma “biasa”
dari $a$ sebagai unsur $k$-aljabar $A$.
Maka kita mempunyai
\[
\chi_{A,a}(t)=\operatorname{Prd}_{A,a}^n(t) \quad
\operatorname{Tr}_{A/k}(a)=n\operatorname{Trd}_A(a) \quad
\operatorname{N}_{A/k}(a)=\left(\operatorname{Nrd}_A(a)\right)^n
\]
dengan $n=\deg(A)$.
\end{proposition}

\subsection{Aljabar Kuaternion}

Misalkan $k$ suatu medan berkarakteristik $\ne 2$.
Andaikan $a,b \in k^\times$.
Misalkan $Q_{a,b}$ suatu $k$-aljabar dengan basis $\{1,i,j,ij\}$ yang
perkaliannya memenuhi
\[
i^2=a, j^2=b, ij=-ij .
\]
Aljabar $Q_{a,b}$ adalah \emph{$k$-aljabar kuaternion (tergeneralisasi)}.
Amati bahwa $Q_{-1,-1} \cong \mathbb{H}$ ketika $k=\mathbb{R}$.

Misalkan $K=k(\sqrt{a})$. Kita mempunyai homomorfisme $k$-aljabar
injektif $\Phi : Q_{a,b} \to \operatorname{M}_2(K)$ melalui
\[
x+yi+zj+wij \mapsto
\begin{pmatrix}
x+\sqrt{a}y & bz+b\sqrt{a}w\\
z-\sqrt{a}w & x-\sqrt{a}y
\end{pmatrix}
\]
untuk $x,y,z,w \in k$.
Ini menunjukkan bahwa $Q_{a,b}$ merupakan $k$-aljabar sederhana sentral dan $K$
suatu medan pengurai.
Perhatikan bahwa $k(\sqrt{b})$ juga merupakan medan pengurai dari $Q_{a,b}$.

Norma tereduksi diberikan oleh
\[
\operatorname{Nrd}(x+yi+zj+wij) = x^2-ay^2-bz^2+abw^2 .
\]
Amati bahwa jika $q \in Q_{a,b}$, maka $\Phi(q)$ invertibel
jika dan hanya jika $\operatorname{Nrd}(q) \ne 0$.
Dalam kasus ini, invers dari $q$ adalah
$\overline{q}/\operatorname{Nrd}(q)$, dengan
\[
\overline{q} = x-yi-zj-wij
\]
merupakan \emph{konjugat} dari $q$.
Jadi, invers dari $q$, jika ada dalam $\operatorname{M}_2(K)$,
juga berada dalam $Q_{a,b}$.

\begin{proposition}
Misalkan $Q_{a,b}$ suatu $k$-aljabar kuaternion.
Pernyataan-pernyataan berikut ekuivalen:
\begin{enumerate}
\item[(a)] $Q_{a,b}$ terurai.
\item[(b)] $Q_{a,b}$ bukan $k$-aljabar divisi.
\item[(c)] $\operatorname{Nrd} : Q_{a,b} \to k$ mempunyai nol nontrivial. 
\item[(d)] $ax^2+by^2=1$ mempunyai solusi $(x,y)$ atas $k$.
\end{enumerate}
\end{proposition}

\begin{proof}
Ekuivalensi (a) dan (b) diperoleh dari teorema Wedderburn
karena $Q_{a,b}=M_n(D)$ untuk suatu $k$-aljabar divisi $D$, dan
$n$ hanya mungkin bernilai $1$ atau $2$.
Ekuivalensi (b) dan (c) diperoleh dari pembahasan invers
di atas.
Kita melihat bahwa (d) mengakibatkan (c) melalui unsur $1-xi-yj$.

Sekarang andaikan (c).
Jika $a$ suatu kuadrat, pembuktian sudah selesai. Jadi,
kita dapat mengasumsikan $k(\sqrt{a})/k$
merupakan perluasan medan nontrivial.
Maka terdapat $x,y,z,w \in k$, tidak semuanya nol,
sedemikian sehingga
\[
(x^2-ay^2)=b(z^2-aw^2).
\]
Amati bahwa $N(x+y\sqrt{a})=x^2-ay^2$, dengan $N : k(\sqrt{a}) \to k$
norma biasa dari perluasan medan $k(\sqrt{a})/k$.
Jadi, $b=N(x+y\sqrt{a})N(z+w\sqrt{a})^{-1}=N(u+v\sqrt{a})$
untuk suatu $u,v \in k$ berdasarkan kemultiplikatifan peta norma.
Dengan kata lain, $b = u^2-av^2$, sehingga kita mempunyai $a(v/u)^2 + b(1/v)^2=1$.
Jadi, $(u/v,1/v)$ memberikan solusi yang diinginkan dan menegakkan (d).
\end{proof}

\subsection{Aljabar Siklik}

Misalkan $L/k$ suatu perluasan Galois siklik berderajat $n$,
andaikan $\sigma$ suatu pembangkit $\operatorname{Gal}(K/k)$,
dan andaikan $b \in k$.
Perhatikan matriks-matriks
\[
Y = \begin{pmatrix}
0 & 0 & \cdots & 0 & b\\
1 & 0 & \cdots & 0 & 0\\
0 & 1 & \cdots & 0 & 0\\
\vdots & \vdots & \vdots &\ddots & \vdots\\
0 & 0 & \cdots & 1 & 0
\end{pmatrix}
\quad
X(\ell) = \begin{pmatrix}
\ell & 0 & 0 &\cdots & 0\\
0 & \sigma(\ell) & 0 &\cdots & 0\\
0 & 0 & \sigma^2(\ell) &\cdots & 0\\
\vdots & \vdots & \vdots &\ddots & \vdots\\
0 & 0 & 0 &\cdots & \sigma^{n-1}(\ell)
\end{pmatrix}
\]
dengan $\ell \in L$.
Amati bahwa $YX(\ell)=X(\sigma(\ell))Y$, $Y^n=1$,
dan $X(\ell)$ memenuhi polinomial minimal yang sama dengan $\ell$.

Perhatikan $k$-aljabar $A=A(L,\sigma,b)$ yang dibangkitkan oleh $X$ dan $Y$.
Aljabar semacam itu disebut \emph{aljabar siklik}.
Amati bahwa, setidaknya ketika $a$ bukan kuadrat dalam $k$,
aljabar kuaternion di atas merupakan kasus khusus
dengan $Q_{a,b} \cong A(k(\sqrt{a}),\sqrt{a}\mapsto -\sqrt{a},b)$.

Pilih $X=X(\ell)$ dengan $\ell$ suatu unsur primitif dari $L/k$.
Amati bahwa $X^n$ berada dalam rentang himpunan bebas linear atas $k$
$\{1,X,X^2,\ldots,X^{n-1}\}$.
Kita melihat bahwa
\[
\{ X^iY^j \mid 0 \le i,j < n \}
\]
merupakan basis bagi $A$ sebagai ruang vektor atas $k$.
Karena $\dim_k A = n^2$ dan kita mempunyai pembenaman sebagai $k$-aljabar
ke dalam $M_n(L)$, kita melihat bahwa setiap aljabar siklik merupakan aljabar sederhana
sentral berderajat $n$.

\subsection{Grup Brauer atas medan khusus}

Secara umum, aljabar sederhana sentral atas medan sebarang dapat sangat
sulit dipahami. Bidang ini benar-benar merupakan ranah penelitian yang
aktif dan terus berkembang. Namun, dalam kasus-kasus khusus yang sangat relevan
bagi teori representasi grup hingga, kita dapat mengatakan jauh
lebih banyak. Sekarang kita merangkum beberapa di antaranya.

Pertama-tama, sebuah pengamatan yang sepele tetapi penting:

\begin{proposition}
Jika $k$ tertutup secara aljabar, maka $\operatorname{Br}(k)=0$.
Dengan kata lain, setiap aljabar sederhana sentral atas medan hingga terurai. 
\end{proposition}

Berikutnya, kita meninjau bilangan real.

\begin{theorem}[Frobenius]
Satu-satunya $\mathbb{R}$-aljabar divisi sentral adalah
$\mathbb{R}$ dan $\mathbb{H}$.
Khususnya, $\operatorname{Br}(\mathbb{R})=\mathbb{Z}/2\mathbb{Z}$
\end{theorem}

Demikian pula untuk medan hingga:

\begin{theorem}[Teorema Kecil Wedderburn]
Jika $k$ medan hingga, maka $\operatorname{Br}(k)=0$.
Dengan kata lain, setiap aljabar sederhana sentral atas medan hingga terurai.
\end{theorem}

Ingat bahwa \emph{medan bilangan} $k$ adalah perluasan hingga dari
$\mathbb{Q}$. Secara kasar, segala hal menarik dalam
teori representasi grup hingga terjadi atas medan bilangan.
Berkat karya Albert--Brauer--Hasse--Noether, kita mempunyai jawaban yang sangat lengkap
dalam kasus ini. Yakni, konstruksi dalam bagian sebelumnya
cukup umum untuk mencakup setiap contoh:

\begin{theorem}
Misalkan $k$ suatu medan bilangan.
Setiap $k$-aljabar sederhana sentral ekuivalen Brauer dengan suatu
$k$-aljabar siklik.
\end{theorem}

Sebenarnya, kita mengetahui lebih banyak lagi jika mengenal sedikit teori bilangan aljabar.
Ingat bahwa medan bilangan $k$ mempunyai \emph{prima-prima}
yang mungkin “hingga” atau “tak hingga”.
Untuk setiap prima, terdapat medan lokal $k_p$ yang bersesuaian.
Contoh prototipe ialah $k=\mathbb{Q}$, dengan prima-prima $p$ berupa
prima biasa dalam $\mathbb{N}$ beserta prima real “tak hingga”
$\infty$. Di sini $\mathbb{Q}_p$ merupakan medan $p$-adik untuk setiap
prima hingga $p$, dan $\mathbb{Q}_\infty = \mathbb{R}$.

\begin{theorem}
Misalkan $k$ medan bilangan dengan prima $p$. Maka
\[
\operatorname{Br}(k_p) \cong
\begin{cases}
\mathbb{Q}/\mathbb{Z} & \textrm{jika $p$ hingga,}\\
\frac{1}{2}\mathbb{Z}/\mathbb{Z} & \textrm{jika $p$ real,}\\
0 & \textrm{jika $p$ kompleks.}\\
\end{cases}
\]
\end{theorem}

Khususnya, teorema di atas memungkinkan kita mendefinisikan
\emph{invarian Hasse}:
\[
\operatorname{inv}_p : \operatorname{Br}(k_p) \to \mathbb{Q}/\mathbb{Z}
\]
yang merupakan homomorfisme grup injektif.
Notasi $\frac{1}{2}\mathbb{Z}/\mathbb{Z}$ untuk kasus real
dimaksudkan untuk menekankan bahwa $\operatorname{inv}_\infty([\mathbb{R}])=1$,
sedangkan $\operatorname{inv}_\infty([\mathbb{H}])=\frac{1}{2}$.

Diberikan aljabar divisi sentral $A$ atas medan bilangan $k$
dan prima $p$, perluasan dasar $A \otimes_k k_p$ merupakan
$k_p$-aljabar divisi sentral.
Jadi, untuk setiap prima $p$, kita mempunyai homomorfisme grup
\[
\operatorname{Br}(k) \to \operatorname{Br}(k_p).
\]
Sekarang kita mempunyai prinsip lokal-ke-global berikut,
yang sepenuhnya mengklasifikasikan aljabar divisi sentral atas suatu medan
bilangan.

\begin{theorem}
Misalkan $k$ suatu medan bilangan.
Terdapat barisan eksak grup-grup abelian
\[
0 \to \operatorname{Br}(k) \to
\bigoplus_p \operatorname{Br}(k_p) \to \mathbb{Q}/\mathbb{Z} \to 0
\]
dengan jumlah langsung diambil atas semua prima $p$ dari $k$.
\end{theorem}

\bibliographystyle{alpha}
\bibliography{rep_theory}

\end{document}


